\section{Weighted objectives and many interacting block responses}
\label{sec:main-weighted}

A terminal difference is a weighted potential with two nonzero
coefficients. General weights can create competing local contributions,
so one-active-block localization no longer suffices. This is an actual
complexity issue, not just a limitation of the proof: with growing
objective support, weighted-potential threshold attainment is
NP-complete on degree-three trees with fixed resistances and a balanced
nomination box. The reduction also excludes absolute-error $1/8$
approximation on its unnormalized family. At fixed support $p$, trees
instead admit exact rational extrema, scenarios, flows, and potentials
in $N^{O(p)}$ time (\cref{thm:a-weight-tree-hard,thm:a-weight-tree-algorithm}).

\subsection{Three completed weighted extensions}

\begin{theorem}[Weighted accuracy-bit algorithms]
\label{thm:main-weighted}
Let $W=c^\top\pi$ with $c\in\Q^V$ and $\one^\top c=0$. Consider any of:
fixed positive rational asymmetric quadratic coefficients; independent positive rational
symmetric coefficient intervals; or independent positive rational
directional coefficient intervals. Scenarios have no operating filters.
In each of the following cases, both extrema satisfy the rational
value and original-scenario contract \eqref{eq:main-contract} in time
polynomial in $N_\eps$ for the fixed parameters:
\begin{enumerate}
\item The graph is a cactus. Nominations form a rational balanced affine
family on a nonempty compact rational polytope of fixed dimension $d$; objective
support and total cycle rank are unrestricted. Alternatively, nominations
range over a balanced rational box and $p=|\supp c|$ is fixed.
\item Maximum block rank is at most a fixed $r_0$, and nominations
range on one nonempty compact rational affine line segment. Objective
support, number of noncactus blocks, and total cycle rank are unrestricted.
\item Put $\kappa(G)=\sum_{B:r_B\ge2}r_B$. The total noncactus rank
$\kappa(G)$ is bounded by a fixed $\kappa_0$, and the balanced affine
nomination domain has fixed dimension $d$. Alternatively, use a balanced
rational nomination box and fixed objective support $p$. Cactus rank
is unrestricted.
\end{enumerate}
In case 2 with fixed laws, the conclusion also holds for admissible
densely encoded rational continuous piecewise-polynomial laws vanishing
at zero, with fixed rational breakpoints and growing degrees and piece
counts. That extension returns rational nominations and does not assert
uncertain polynomial coefficients.
\end{theorem}

Case 1 is \cref{thm:a-wcac-affine,cor:a-wcac-box}; case 2 is
\cref{thm:a-wblk-line,cor:a-wblk-line-polynomial}; case 3 is
\cref{thm:a-wblk-hybrid,cor:a-wblk-hybrid-box}.
For comparison, fixed \emph{global} cycle rank $r$ and fixed support
$p$ permit exact algebraic weighted extrema and algebraic scenarios
under independent continuous resistance intervals
(\cref{thm:a-weight-global}). There all nonlinear variables fit in one
common algebraic field. Fixed rank per block allows an unbounded number
of independent fields, so it requires the separate summation arguments.

The graph and support parameters can sometimes be reduced further.
On a tree, contract every edge with zero induced objective cut flow,
then orient each remaining edge along positive cut flow. Only vertices
that fail to have exactly one incoming and one outgoing edge need be
marked. If their number $k$ is fixed, exact optimization takes $N^{O(k)}$
time even with growing support. At fixed global rank a supplied marked
set and constant source sign on each suppressed path give a related
signed-path refinement, including returning paths. These statements
are formalized in \cref{thm:a-weight-tree-algorithm,thm:a-weight-global};
they describe checkable structural information, not an assumption that
all objective weights have the same sign.

\subsection{A face theorem and two ways to sum}

Choose a rational tree flow $w$ satisfying $Aw=c$. Then
\[
 W=c^\top\pi=w^\top A^\top\pi=\sum_e w_eg_e(x_e).
\]
This puts each contribution in its physical block. Prune subtrees with
no objective support and contract a block if every induced objective
sum at its attachments vanishes. Both operations preserve the objective
and admit rational nomination disaggregation. The incidence subtree
spanning $p$ objective vertices has only $O(p)$ branching locations
and ordinary chains. An adjoint threshold can leave one selected block
per ordinary chain plus the special blocks. Path suppression and a
second positive-source perturbation then leave $O(r_0p)$ free nominations
on $n^{O(r_0p)}$ closed faces for a positive bound $r_0\ge1$ with
$r_{\max}\le r_0$ (\cref{thm:a-wcac-faces}). This full weighted
face reduction is independent of the coefficients. It is structural;
it does not imply an additive algorithm for an arbitrary number of
higher-block value functions of several parameters.

For a quadratic cactus, a local circulation equation is piecewise
quadratic. With uncertain coefficients, a single equality LP first
eliminates the resistance choices. Reduced costs are ordered by one
weight threshold; tied weights remain as a recoverable interval sum.
The complete circulation candidates include boundaries, stationary
points, flat pieces, and zero-flow cases. A local maximum is therefore
selected from polynomially many rational or quadratic-radical formulas
in the nomination parameters. Geometrically spaced rational polynomial
panels approximate the square roots, and a common sign decomposition
allows all block surrogates to be added in the \emph{same} fixed-dimensional
nomination coordinates. Exact optimization of the surrogate uses a
fixed number of variables; it does not introduce one value variable
for every cycle. Global continuity then recovers original rational
parameters, even if rounding changes local flow signs.

For higher-rank blocks on one nomination line, local algebraic degree
need not be two. The key interface becomes a rational quantifier-free
graph of a scalar Lipschitz semialgebraic function. The construction in
\cref{lem:a-wblk-univariate} isolates exceptional parameters, uses the
Lipschitz bound on small neighborhoods, and constructs rational analytic
interpolants on geometric panels elsewhere. Polynomial dependence on
dense degree, coefficient bits, and accuracy bits is proved explicitly.
The resulting rational polynomials retain the original parameter $t$,
so arbitrarily many block responses can be added and optimized without
forming their joint algebraic field. This is why the one-line result
can allow an unrestricted number of higher-rank blocks.

For several nomination parameters, case 3 instead keeps all noncactus
circulations in one exact core of dimension $d+\kappa(G)$ and compiles
the remaining cactus contributions. The quantity $\kappa(G)$ is a
\emph{sum}, not a maximum. A graph with a thousand rank-two blocks
does not have small $\kappa(G)$ merely because each block has rank two.
The separate one-line theorem does apply to that graph.

\subsection{Exact parameters, filters, and limits}

At fixed rational nominations on a cactus, the quadratic weighted
algorithm gives an exactly optimizing rational coefficient profile,
even with optional signed arc capacities and arbitrary objective support.
Its local values have degree at most two, and the global value may be
retained as their separate sum (\cref{thm:a-wcac-exact-profile}).
An exact profile need not be an endpoint profile: on the four-cycle in
\cref{ex:a-wcac-interior}, the unique optimum is $(1,2,1,1)$ and the
second coefficient lies strictly inside $[1,3]$.
Thus a polynomial exact scenario output is compatible both with an
interior optimum and with the absence of an efficient exact scalar
comparator for its value.

For fixed-dimensional affine nominations, local capacity filters permit
exact feasibility and additive algebraic scenarios of polynomial
representation size. With fixed laws, closed bounded-degree
semialgebraic filters involving the nomination parameters and one block's
flows are also allowed. The joint coefficient theorem permits only the
stated affine arc bounds. A supplied positive tightening margin gives
a rational scenario feasible for the original capacities and near-optimal
relative to the \emph{tightened} optimum
(\cref{thm:a-wcac-capacities}). The reference optimum matters: a small
coefficient perturbation need not preserve an optimizer lying on an
exact capacity boundary.

Indeed rational inputs can force an irrational nomination. The four-cycle
in \cref{ex:a-wcac-irrational} has flows
$(q+t-1,q-2,q,q-2)$ and resistances $(1,1,1,2)$.
The bounds $|x_3|\le1$ and $|x_4|\le1$ force $q=1$; the
cycle equation becomes $t^2=2$ for $1\le t\le2$.
Thus no rational feasible nomination exists. This is an output
obstruction, not a contradiction to the unfiltered rational approximation
theorems.

The unresolved weighted case is now specific: several fixed shared
nomination parameters, an unrestricted number of higher-rank blocks,
and only a fixed \emph{maximum} block rank. The local value graphs and
face reduction are available. What remains outside the proved algorithms
is accuracy-bit optimization of their growing sum in the original
shared coordinates (\cref{eq:a-wblk-residual}). A Lipschitz grid is
polynomial in $1/\eps$, and retaining all local values defeats
fixed-dimensional elimination. These observations delimit those methods;
they are not a hardness or impossibility result.
